\subsection{SVD}

\begin{proposition}
Let $A \in M_n(\mathbb{C})$. Consider the Hermitian matrix $A^*A$.
\begin{enumerate}
    \item $A^*A$ is PSD.
    \item There exists a PSD Hermitian matrix $D$ such that $A^*A=D^2$.
\end{enumerate}
\end{proposition}

\begin{pfbox}
\begin{enumerate}
    \item For any $\mathbf{v}\in \mathbb{C}^n$,
    \[
    \innerproduct{A^*A\mathbf{v}}{\mathbf{v}}
    =
    \innerproduct{A\mathbf{v}}{A\mathbf{v}}
    =
    \norm{A\mathbf{v}}^2
    \ge 0.
    \]
    Thus, all eigenvalues $\lambda\ge 0$, meaning $A^*A$ is PSD.

    \item By the spectral theorem,
    \[
    V^*A^*AV
    =
    \Lambda
    =
    \operatorname{diag}(\lambda_1,\lambda_2,\dots,\lambda_n).
    \]
    Write
    \[
    \Sigma
    =
    \operatorname{diag}(\mu_1,\mu_2,\dots,\mu_n),
    \qquad
    \mu_i=\sqrt{\lambda_i}\ge 0.
    \]
    Then $\Sigma^2=\Lambda$ and $\Sigma^*=\Sigma$. Therefore,
    \[
    A^*A
    =
    V\Lambda V^*
    =
    V\Sigma^2V^*
    =
    (V\Sigma V^*)(V\Sigma V^*)
    =
    D^2.
    \]
    Let
    \[
    D=V\Sigma V^*.
    \]
    Then $D$ is Hermitian and PSD.
\end{enumerate}
\end{pfbox}

\begin{theorem}[SVD: $\mathbb{C}$omplex version]
Let $A \in M_{m \times n}(\mathbb{C})$ be a matrix of rank $k$.
There exist
\begin{enumerate}
    \item unitary matrices $U \in U(m)$ and $V \in U(n)$;
    \item a diagonal matrix
    \[
    \Sigma
    =
    \begin{pmatrix}
    \sigma_1 & 0 & \dots & 0 \\
    0 & \sigma_2 & \dots & 0 \\
    \vdots & \vdots & \ddots & \vdots \\
    0 & 0 & \dots & \sigma_k
    \end{pmatrix}
    \]
    with
    \[
    \sigma_1 \ge \sigma_2 \ge \dots \ge \sigma_k > 0,
    \]
\end{enumerate}
such that $A$ admits a factorization of the form
\[
A
=
U
\underbrace{
\begin{pmatrix}
\Sigma & \mathbf{O} \\
\mathbf{O} & \mathbf{O}
\end{pmatrix}}_{\Sigma_0\in M_{m \times n}(\mathbb{R})}
V^*.
\]
\end{theorem}


\begin{pfbox}
The proof is constructive.
Consider $A\in M_{m\times n}(\mathbb{C})$,
Then $A^*A\in M_n(\mathbb{C})$ is a square matrix.
Let $k=\rank(A)=\rank(A^*A)$ (why?), 
and $\lambda_1\ge \lambda_2\ge \dots\ge \lambda_k>0$
be the non-zero eigenvalues of $A^*A$.\\

\textbf{Step 1. Diagonalize $A^*A$.}
As $A^*A$ is Hermitian, and hence normal, by the Spectral Theorem, there exists
$V\in U(n)$ such that
\[
V^*(A^*A)V
=
\begin{pmatrix}
\Sigma^2 & \mathbf{O} \\
\mathbf{O} & \mathbf{O}
\end{pmatrix}.
\]
Take $\sigma_i=\sqrt{\lambda_i}$, and write $V=(V_1\mid V_2),$
where $V_1$ contains the first $k$ columns.\\

\textbf{Step 2. Compare the blocks of $(AV)^*(AV)$.}
Since
\[
AV=(AV_1\mid AV_2),
\]
compute:
\[
(AV)^*(AV)
=
\begin{pmatrix}
(AV_1)^*(AV_1) & (AV_1)^*(AV_2) \\
(AV_2)^*(AV_1) & (AV_2)^*(AV_2)
\end{pmatrix}
=
\begin{pmatrix}
\Sigma^2 & \mathbf{O} \\
\mathbf{O} & \mathbf{O}
\end{pmatrix}.
\]
By comparing blocks, we have
\[(AV_1)^*(AV_1)=\Sigma^2 = \Sigma^*\Sigma\] 
and
\[(AV_2)^*(AV_2)=\mathbf{O},\] 
which implies $AV_2=\mathbf{O}$ since $\ker((AV_2)^*(AV_2))=\ker((AV_2))$.\\

\textbf{Step 3. Construct $U_1$.}

Let
\[
U_1=AV_1\Sigma^{-1}\in M_{m\times k}(\mathbb{C}).
\]
Note that
\[
U_1^*U_1
=
\Sigma^{-1}V_1^*A^*AV_1\Sigma^{-1}
=
I_k,
\]
so the columns of $U_1$ are orthonormal.
By the Basis Extension Theorem and Gram-Schmidt, we can extend $U_1$ to a unitary
matrix
\[
U=(U_1\mid U_2)\in U(m).
\]

\textbf{Step 4. Verify the decomposition.}

Now compute:
\[
U^*AV
=
\begin{pmatrix}
U_1^*\\
U_2^*
\end{pmatrix}
A(V_1\mid V_2)
=
\begin{pmatrix}
U_1^*AV_1 & U_1^*AV_2 \\
U_2^*AV_1 & U_2^*AV_2
\end{pmatrix}.
\]
Substitute
\[
AV_1=U_1\Sigma,
\qquad
AV_2=\mathbf{O}.
\]
Then
\[
U^*AV
=
\begin{pmatrix}
U_1^*U_1\Sigma & \mathbf{O} \\
U_2^*U_1\Sigma & \mathbf{O}
\end{pmatrix}
=
\begin{pmatrix}
\Sigma & \mathbf{O} \\
\mathbf{O} & \mathbf{O}
\end{pmatrix}.
\]
Thus
\[
A=U\Sigma_0V^*,
\]
completing the SVD.
\end{pfbox}

\begin{remark}[Vector form of SVD]
Let $\{\mathbf{u}_1,\dots,\mathbf{u}_m\}$
be the columns of $U$ and
$\{\mathbf{v}_1,\dots,\mathbf{v}_n\}$
be the columns of $V$. 
The vector form of the SVD of $A$ is the sum of rank-one matrices:
\[
A=\sum_{i=1}^k \sigma_i\mathbf{u}_i\mathbf{v}_i^*.
\]
\end{remark}



\begin{theorem}[Operator Form of SVD]
Let
\[
T:W_1\rightarrow W_2
\]
be a linear map between two finite-dimensional complex inner product spaces.
There exist
\begin{enumerate}
    \item an ONB $\{\mathbf{v}_1,\dots,\mathbf{v}_n\}$ of $W_1$ and an ONB
    $\{\mathbf{u}_1,\dots,\mathbf{u}_m\}$ of $W_2$;
    \item positive values
    \[
    \sigma_1\ge \sigma_2\ge \dots\ge \sigma_k>0,
    \]
\end{enumerate}
such that, where $k=\rank(T)$,
\[
T(\mathbf{v}_i)
=
\begin{cases}
\sigma_i \mathbf{u}_i, & \text{if } i\le k,\\
\mathbf{0}, & \text{otherwise.}
\end{cases}
\]
\end{theorem}


\begin{definition}[Singular values]
Let $A\in M_{m\times n}(\mathbb{C})$.
The values
\[
\sigma_1\ge \sigma_2\ge \dots\ge \sigma_k>0
\]
that occur in the SVD of $A$ are called the singular values of $A$.
\end{definition}

\begin{theorem}
If $A$ admits an SVD
\[
A=U\Sigma_0V^*,
\]
then the diagonal entries of $\Sigma$,
\[
\sigma_1\ge \sigma_2\ge \dots\ge \sigma_k>0,
\]
are given by the positive square roots of the non-zero eigenvalues of the Hermitian
matrix $A^*A$.
In particular, the values $\sigma_i$ are uniquely determined by $A$.
\end{theorem}

\begin{pfbox}
Suppose $A$ admits an SVD of the form
\[
A=U\Sigma_0V^*.
\]
Then we compute $A^*A$:
\begin{align*}
A^*A
&=
\left(
U
\begin{pmatrix}
\Sigma & \mathbf{O} \\
\mathbf{O} & \mathbf{O}
\end{pmatrix}
V^*
\right)^*
\left(
U
\begin{pmatrix}
\Sigma & \mathbf{O} \\
\mathbf{O} & \mathbf{O}
\end{pmatrix}
V^*
\right) \\
&=
V
\begin{pmatrix}
\Sigma^* & \mathbf{O} \\
\mathbf{O} & \mathbf{O}
\end{pmatrix}
U^*U
\begin{pmatrix}
\Sigma & \mathbf{O} \\
\mathbf{O} & \mathbf{O}
\end{pmatrix}
V^* \\
&=
V
\begin{pmatrix}
\Sigma^2 & \mathbf{O} \\
\mathbf{O} & \mathbf{O}
\end{pmatrix}
V^*.
\end{align*}
Since $V$ is unitary, this means $A^*A$ is unitarily similar to a diagonal matrix
with entries $\sigma_i^2$.
So $\sigma_i^2$ is an eigenvalue of $A^*A$, and $\sigma_i$ is the positive square
root of the eigenvalue of $A^*A$.
\end{pfbox}


\begin{theorem}[Polar Decomposition]
Let $A \in M_n(\mathbb{C})$. Then there exists a unitary matrix $L \in U(n)$ and a positive semi-definite (PSD) Hermitian matrix $P$ such that
\[
A=LP.
\]
Moreover, if $A$ is invertible, then the factorization is unique.
\end{theorem}

\begin{pfbox}
\textbf{Existence.}
By SVD,
\[
A=U\Sigma V^*,
\]
where $U,V\in U(n)$ and $\Sigma$ is PSD diagonal. Hence
\[
A=(UV^*)(V\Sigma V^*)=LP.
\]
Here,
\[
L=UV^*,
\qquad
P=V\Sigma V^*.
\]
Then $L$ is unitary, and $P$ is Hermitian and PSD.

\medskip

\textbf{Uniqueness.}
Suppose there exist $L_1,L_2$ and $P_1,P_2$ such that
\[
A=L_1P_1=L_2P_2.
\]
Then
\[
A^*A=P_1^*L_1^*L_1P_1=P_1^*P_1=P_1^2
\]
and
\[
A^*A=P_2^*L_2^*L_2P_2=P_2^*P_2=P_2^2.
\]
Since $P_1,P_2$ are Hermitian and PSD (actually PD if $A$ is invertible), if
$P_1^2=P_2^2$, then $P_1=P_2$ (Exercise), and hence $L_1=L_2$.
\end{pfbox}

\subsubsection*{Applications: Low-Rank Approximation}

\begin{theorem}[Von Neumann's Trace Inequality]
Let $A,B\in M_{m\times n}(\mathbb{R})$. Suppose $\rank(A)=k$ and $\rank(B)=k'$,
and let $\sigma_1\ge\sigma_2\ge\cdots\ge\sigma_k>0$ and $\mu_1\ge\mu_2\ge\cdots\ge\mu_{k'}>0$ be the singular values of $A$ and $B$, respectively. Then:
\[|\tr(AB^\top)|\le\sum_{i=1}^{\min\{k,k'\}}\sigma_i\mu_i.\]
\end{theorem}

\begin{pfbox}
Define $\displaystyle\sum_{i=1}^{\min\{k,k'\}}\sigma_i\mu_i = \sum_{i=1}^{N}\sigma_i\mu_i,$ where $N>k,k'$.\\
Set $\sigma_i=0$ if $i>k$, and $\mu_i=0$ if $i>k'$.
By SVD,
\[
    A = \sum_{i=1}^{N}\sigma_i\mathbf{u}_i\mathbf{v}_i^\top,\quad
    B = \sum_{i=1}^{N}\sigma_i\mathbf{x}_i\mathbf{y}_i^\top,
\]

where $\{\mathbf{u}_i\}^m_{i=1}$ and $\{\mathbf{x}_i\}^m_{i=1}$ are ONBs of $\mathbb{R}^m$, and $\{\mathbf{v}_i\}^n_{i=1}$ and $\{\mathbf{y}_i\}^n_{i=1}$ are ONBs of $\mathbb{R}^n$.

Define, for $\ell,s\le N,\, A_\ell = \displaystyle\sum_{i=1}^{\ell}\mathbf{u}_i\mathbf{v}_i^\top,\,\, B_s = \displaystyle\sum_{j=1}^{s}\mathbf{x}_j\mathbf{y}_j^\top$.\\
Von Neumann \textcolor{red}{observes} $A=\displaystyle\sum_{\ell=1}^{N}(\sigma_\ell-\sigma_{\ell+1})A_\ell$ and $B = \displaystyle\sum_{s=1}^{N}(\mu_s-\mu_{s+1})B_s$. 

Then $\tr(AB^\top)=\displaystyle\sum_{\ell=1}^{N}\displaystyle\sum_{s=1}^{N}(\sigma_\ell-\sigma_{\ell+1})(\mu_s-\mu_{s+1})\tr(A_\ell B_s^\top).$

Next, we show that $\tr(A_\ell B_s^\top)\le\min\{\ell,s\}$ for any $\ell,s\le N.$
\begin{align*}
    |\tr(A_\ell B_s^\top)| &= |\sum_{i=1}^{\ell}\sum_{j=1}^{s}\tr(\mathbf{u}_i\mathbf{v}_i^\top\mathbf{y}_j\mathbf{x}_j^\top)|\\
    &= |\sum_{i=1}^{\ell}\sum_{j=1}^{s}\tr(\mathbf{v}_i^\top\mathbf{y}_j\mathbf{x}_j^\top \mathbf{u}_i)| = |\sum_{i=1}^{\ell}\sum_{j=1}^{s}\mathbf{v}_i^\top\mathbf{y}_j\mathbf{x}_j^\top \mathbf{u}_i|\\
    &= |\sum_{i=1}^{\ell}\sum_{j=1}^{s} \langle\mathbf{y}_j\mathbf{x}_j^\top \mathbf{u}_i, \mathbf{v}_i\rangle| =  |\sum_{i=1}^{\ell}\sum_{j=1}^{s} \langle \mathbf{v}_i, \mathbf{y}_j\mathbf{x}_j^\top \mathbf{u}_i\rangle|.
\end{align*}

Let $\mathbf{w}_i=\displaystyle\sum_{j=1}^{s}\mathbf{y}_j\mathbf{x}_j^\top \mathbf{u}_i=\displaystyle\sum_{j=1}^{s}\langle \mathbf{x}_j, \mathbf{u}_i\rangle \mathbf{y}_j.$
Claim: $\norm{\mathbf{w}_i} = 1$ for any $i\in\{1,2,\cdots,\ell\}.$
By Pythagorean theorem,
\begin{align*}
\norm{\mathbf{w}_i}&=\sqrt{\sum_{j=1}^{s}\norm{\langle \mathbf{x}_j, \mathbf{u}_i\rangle \mathbf{y}_j}^2}
\le\sqrt{\sum_{j=1}^{m}\norm{\langle \mathbf{x}_j, \mathbf{u}_i\rangle \mathbf{y}_j}^2}
=\sqrt{\sum_{j=1}^{m}|\langle \mathbf{x}_j, \mathbf{u}_i\rangle |^2\norm{\mathbf{y}_j}^2}\\
&=\sqrt{\sum_{j=1}^{m}|\langle \mathbf{x}_j, \mathbf{u}_i\rangle |^2} 
=\norm{\mathbf{u}_i}=1.  \quad \text{($\{y_j\}$ and $\{x_j\}$ are ONBs)}
\end{align*}

So
\begin{align*}
|\tr(A_\ell B_s^\top)|& = |\sum_{i=1}^{\ell}\sum_{j=1}^{s} \langle \mathbf{v}_i, \mathbf{y}_j\mathbf{x}_j^\top \mathbf{u}_i\rangle|\\
& = |\sum_{i=1}^{\ell}\langle \mathbf{v}_i, w_i\rangle|\le \sum_{i=1}^{\ell}\norm{\mathbf{v}_i}\norm{\mathbf{w}_i}=\sum_{i=1}^{\ell}1\times 1 = \ell.    
\end{align*}

By symmetry, we also have $|\tr(A_\ell B_s^\top)|\le s$.

Hence,
\begin{align*}
|\tr(AB^\top)|&=|\sum_{\ell=1}^{N}\sum_{s=1}^{N}(\sigma_\ell-\sigma_{\ell+1})(\mu_s-\mu_{s+1})\tr(A_\ell B_s^\top)|\\
&\le \sum_{\ell=1}^{N}\sum_{s=1}^{N}(\sigma_\ell-\sigma_{\ell+1})(\mu_s-\mu_{s+1})|\tr(A_\ell B^\top_s)|\\
&\le\sum_{\ell=1}^{N}\sum_{s=1}^{N}(\sigma_\ell-\sigma_{\ell+1})(\mu_s-\mu_{s+1})\min\{\ell,s\}\\
&\le \sum_{s=1}^{N}(\mu_s-\mu_{s+1})(\sum_{\ell=1}^{s}(\sigma_\ell-\sigma_{\ell+1})\ell + \sum_{\ell=s+1}^{N}(\sigma_\ell-\sigma_{\ell+1})s)\\
&\le \sum_{s=1}^{N}(\mu_s-\mu_{s+1})(\sum_{\ell=1}^{s}\sigma_\ell - s\sigma_{s+1} + s\sigma_{s+1})\\
&\le \sum_{s=1}^{N}(\mu_s-\mu_{s+1})\sum_{\ell=1}^{s}\sigma_\ell = \sum_{s=1}^{N}\mu_s\sigma_s = \sum_{i=1}^{\min\{k,k'\}}\sigma_i\mu_i.
\end{align*}
This completes the proof.
\end{pfbox}



\begin{theorem}[Low-rank approximation]
Let $A\in M_{m\times n}(\mathbb{R})$ be of rank $k$ with singular values
\[
\sigma_1\ge \sigma_2\ge \dots\ge \sigma_k>0.
\]
Fix a value $l<k$.
For any matrix $B\in M_{m\times n}(\mathbb{R})$ of rank $l$, one has
\[
\|A-B\|_F
\ge
\sqrt{\sigma_{l+1}^2+\sigma_{l+2}^2+\dots+\sigma_k^2}.
\]
The equality holds when
\[
B=
\sum_{i=1}^{l}
\sigma_i\mathbf{u}_i\mathbf{v}_i^t,
\]
which is the truncated SVD of $A$.
\end{theorem}

\begin{pfbox}
Let $\mu_1\ge\mu_2\ge\cdots\mu_l>0$ be the singular values of $B$.
\begin{align*}
    \|A-B\|^2_F &= \tr((A-B)^\top(A-B)) = \tr((A-B)(A-B)^\top)\\
                &= \tr(AA^\top) + \tr(BB^\top) -\tr(AB^\top)-\tr(BA^\top)\\
                &= \sum_{i=1}^{k}\sigma_i^2 + \sum_{i=1}^{l}\mu_i^2-2\tr(AB^\top)\\
                &\ge \sum_{i=1}^{k}\sigma_i^2 + \sum_{i=1}^{l}\mu_i^2-2\sum_{i=1}^{l}\sigma_i\mu_i \quad \text{(Von Neumann's Trace Inequality)}\\
                &=\sum_{i=1}^{l}\sigma_i^2 + \sum_{i=1}^{l}\mu_i^2-2\sum_{i=1}^{l}\sigma_i\mu_i + \sum_{i=l+1}^{k}\sigma_i^2\\
                &=\sum_{i=1}^{l}(\sigma_i-\mu_i)^2 + \sum_{i=l+1}^{k}\sigma_i^2\\
                &\ge \sum_{i=l+1}^{k}\sigma_i^2.
\end{align*}
The equality holds if $\sigma_i=\mu_i$ for $i=1,2,\cdots,l.$
\end{pfbox}



\subsection{Pseudoinverse}

\begin{definition}[Moore--Penrose inverse]
Let $T:V\rightarrow W$ be a linear map between two inner product spaces.
Denote the inverse of
\[
T|_{\Ker(T)^\perp}:
\Ker(T)^\perp
\rightarrow
\Ima(T)
\]
simply by
\[
T^{-1}:
\Ima(T)
\rightarrow
\Ker(T)^\perp.
\]
The pseudoinverse, or Moore--Penrose inverse, of $T$ is a linear map 
$T^\dagger:W\rightarrow V$ defined by
\[
T^\dagger(w)
=
T^{-1}\bigl(P_{\Ima(T)}(w)\bigr).
\]
\end{definition}

\begin{theorem}
Let $T:V\rightarrow W$ be a linear map between two inner product spaces.
Then its pseudoinverse $T^\dagger$ satisfies:
\begin{enumerate}
    \item $T^\dagger T:V\rightarrow V$ coincides with
    $P_{\Ker(T)^\perp}$;
    \item $TT^\dagger:W\rightarrow W$ coincides with
    $P_{\Ima(T)}$.
\end{enumerate}
\end{theorem}

\begin{pfbox}
By orthogonal decomposition, $V=\Ker(T)\oplus \Ker(T)^\perp$.

Let $\mathbf{v}\in V$, $\mathbf{v}=\mathbf{v}_1+\mathbf{v}_2$,
where $\mathbf{v}_1\in \Ker(T),\,\mathbf{v}_2\in \Ker(T)^\perp$.

Thus
\begin{align*}
T^\dagger T(\mathbf{v}) &= T^\dagger(T(\mathbf{v}_2)) \\
&=T^{-1}\bigl(P_{\Ima(T)}(T(\mathbf{v}_2))\bigr) = T^{-1}(T(\mathbf{v}_2))\\
&= \mathbf{v}_2 = P_{\Ker(T)^\perp}(\mathbf{v}).
\end{align*}

Similarly, for $W=\Ima(T)\oplus \Ima(T)^\perp$, 
let $\mathbf{w}=\mathbf{w}_1+\mathbf{w}_2$.
Then
\begin{align*}
TT^\dagger(\mathbf{w})
&= T\Bigl(T^{-1}\bigl(P_{\Ima(T)}(\mathbf{w}_1+\mathbf{w}_2)\bigr)\Bigr) \\
&=T(T^{-1}(\mathbf{w}_1)) = \mathbf{w}_1 = P_{\Ima(T)}(\mathbf{w}).
\end{align*}
\end{pfbox}


\begin{definition}[Computational definition of pseudoinverse]
Let $A\in M_{m\times n}(\mathbb{C})$ have an SVD
\[
A=
U
\begin{pmatrix}
\Sigma & \mathbf{O} \\
\mathbf{O} & \mathbf{O}
\end{pmatrix}
V^*.
\]
Then we define its pseudo-inverse to be
\[
A^\dagger
\coloneqq
V
\begin{pmatrix}
\Sigma^{-1} & \mathbf{O} \\
\mathbf{O} & \mathbf{O}
\end{pmatrix}
U^*
\in
M_{n\times m}(\mathbb{C}).
\]
\end{definition}

\begin{theorem}[Penrose]
Let $A\in M_{m\times n}(\mathbb{C})$.
Then its pseudoinverse satisfies the Penrose conditions:
\begin{enumerate}
    \item $AA^\dagger A=A$.
    \item $A^\dagger AA^\dagger=A^\dagger$.
    \item $(AA^\dagger)^*=AA^\dagger$.
    \item $(A^\dagger A)^*=A^\dagger A$.
\end{enumerate}
\end{theorem}

\begin{pfbox}
Substitute the SVD forms
\[
A
=
U\Sigma_0V^*,
\qquad
A^\dagger
=
V\Sigma_0^{-1}U^*,
\]
where
\[
\Sigma_0
=
\begin{pmatrix}
\Sigma & \mathbf{O} \\
\mathbf{O} & \mathbf{O}
\end{pmatrix},
\qquad
\Sigma_0^{-1}
=
\begin{pmatrix}
\Sigma^{-1} & \mathbf{O} \\
\mathbf{O} & \mathbf{O}
\end{pmatrix},
\]
and
\[
\Sigma^{-1}
=
\operatorname{diag}
(\sigma_1^{-1},\ldots,\sigma_k^{-1}).
\]

\textbf{Condition 1.}
\begin{align*}
AA^\dagger A
&=
U\Sigma_0V^*
V\Sigma_0^{-1}U^*
U\Sigma_0V^* \\
&=
U\Sigma_0\Sigma_0^{-1}\Sigma_0V^* \\
&=
U\Sigma_0V^* \\
&=
A.
\end{align*}

\textbf{Condition 2.}
\begin{align*}
A^\dagger AA^\dagger
&=
V\Sigma_0^{-1}U^*
U\Sigma_0V^*
V\Sigma_0^{-1}U^* \\
&=
V\Sigma_0^{-1}\Sigma_0\Sigma_0^{-1}U^* \\
&=
V\Sigma_0^{-1}U^* \\
&=
A^\dagger.
\end{align*}

\textbf{Condition 3.}
Since
\begin{align*}
AA^\dagger
&=
U\Sigma_0V^*
V\Sigma_0^{-1}U^* \\
&=
U\Sigma_0\Sigma_0^{-1}U^* \\
&=
U
\begin{pmatrix}
I_k & \mathbf{O} \\
\mathbf{O} & \mathbf{O}
\end{pmatrix}
U^*,
\end{align*}
and
\[
\begin{pmatrix}
I_k & \mathbf{O} \\
\mathbf{O} & \mathbf{O}
\end{pmatrix}
\]
is Hermitian, we get
\[
(AA^\dagger)^*
=
AA^\dagger.
\]

\textbf{Condition 4.}
Since
\begin{align*}
A^\dagger A
&=
V\Sigma_0^{-1}U^*
U\Sigma_0V^* \\
&=
V\Sigma_0^{-1}\Sigma_0V^* \\
&=
V
\begin{pmatrix}
I_k & \mathbf{O} \\
\mathbf{O} & \mathbf{O}
\end{pmatrix}
V^*,
\end{align*}
and
\[
\begin{pmatrix}
I_k & \mathbf{O} \\
\mathbf{O} & \mathbf{O}
\end{pmatrix}
\]
is Hermitian, we get
\[
(A^\dagger A)^*
=
A^\dagger A.
\]

Therefore, $A^\dagger$ satisfies all four Penrose conditions.
\end{pfbox}

\begin{theorem}[Penrose Uniqueness]
If $X\in M_{n\times m}(\mathbb{C})$ satisfies the four Penrose conditions, then
\[
X=A^\dagger.
\]
\end{theorem}

\begin{pfbox}
Suppose both $X_1$ and $X_2$ satisfy the Penrose conditions.
First observe the following two facts.

\textbf{Observation 1.}
By $(AX)^*=AX$ and $XAX=X$, we have
\[
X(AX)^*=XAX=X.
\]
Since $(AX)^*=X^*A^*$, this gives
\[
XX^*A^*=X.
\]

\textbf{Observation 2.}
By $AXA=A$ and $(XA)^*=XA$, we have
\[
A(XA)^*=AXA=A.
\]
Since $(XA)^*=A^*X^*$, this gives
\[
AA^*X^*=A.
\]
Taking adjoints gives
\[
XA A^*=A^*.
\]

Now 
\begin{align*}
X_1
&= X_1X_1^*A^*
    && (\text{Obs 1 for } X_1) \\[2mm]
&= X_1X_1^*(AX_2A)^*
    && (\text{Condition 1 for } X_2) \\[2mm]
&= X_1X_1^*A^*(AX_2)^* \\[2mm]
&= X_1X_1^*A^*(AX_2)
    && (\text{Condition 3 for } X_2) \\[2mm]
&= X_1\textcolor{red}{A}X_2
    && (\text{Obs 1 for } X_1) \\[2mm]
&= X_1\textcolor{red}{AA^*X^*_2}X_2
    && (\text{Obs 2 for } X_2) \\[2mm]
&= A^*X^*_2X_2
  && (\text{Obs 2 for } X_1) \\[2mm]
&= (X_2A)^*X_2 = X_2AX_2 = X_2
  && (\text{Condition 4 \& 2 for } X_2). \\[2mm]
\end{align*}
\end{pfbox}

\begin{theorem}
    Theoretical definition and computational definition of pseudoinverse are equivalent.
\end{theorem}

\begin{pfbox}
Let $T^\dagger:W\to V$ be the pseudoinverse of $T:V\to W$.
Recall $T^\dagger T=P_{\ker(T)^\perp}$ and $TT^\dagger=P_{\Ima(T)}$.

For $v\in V$, we have 
$TT^\dagger T(v)=P_{\Ima(T)}(T(v))=T(v)$. Hence
\[
TT^\dagger T=T.
\]

For $w\in W$, we have
$T^\dagger TT^\dagger(w)
=P_{\Ker(T)^\perp}(T^\dagger(w))
=T^\dagger(w)$,\\
since $T^\dagger(w)\in\ker(T)^\perp$. Hence
\[
T^\dagger TT^\dagger=T^\dagger.
\]


Let $\beta=\{v_1,\dots,v_k,v_{k+1},\dots,v_n\}$ be an ONB of $V$, where
\[
\spn\{v_1,\dots,v_k\}=\ker(T)^\perp,
\qquad
\spn\{v_{k+1},\dots,v_n\}=\ker(T).
\]
Then
\[
\left[P_{\ker(T)^\perp}\right]_{\beta}^{\beta}
=
\begin{pmatrix}
I_k & \mathbf{O}\\
\mathbf{O} & \mathbf{O}_{n-k}
\end{pmatrix},
\]
which is Hermitian.

Similarly, let $\gamma=\{w_1,\dots,w_k,w_{k+1},\dots,w_m\}$
be an ONB of $W$, where
\[
\spn\{w_1,\dots,w_k\}=\Ima(T),
\qquad
\spn\{w_{k+1},\dots,w_m\}=\Ima(T)^\perp.
\]
Then
\[
\left[P_{\Ima(T)}\right]_{\gamma}^{\gamma}
=
\begin{pmatrix}
I_k & \mathbf{O}\\
\mathbf{O} & \mathbf{O}_{m-k}
\end{pmatrix},
\]
which is Hermitian.

Since orthogonal projections are Hermitian, we get
\[
(T^\dagger T)^*=T^\dagger T,
\qquad
(TT^\dagger)^*=TT^\dagger.
\]

Let $[T]_{\beta}^{\gamma}=A$, by Penrose conditions, $[T^\dagger]_{\gamma}^{\beta}=A^\dagger$.
\end{pfbox}

\subsubsection*{Solve System of Equations}

\begin{theorem}[Least Square Method]
Let $A\in M_{m\times n}(\mathbb{C})$ and fix $\mathbf{b}\in \mathbb{C}^m$.
Then $\mathbf{x}=A^\dagger\mathbf{b}$ minimizes the norm
$\|A\mathbf{x}-\mathbf{b}\|$.
That is,
\[
\|A(A^\dagger\mathbf{b})-\mathbf{b}\|
\le
\|A\mathbf{x}-\mathbf{b}\|
\qquad
\text{for all } \mathbf{x}\in \mathbb{C}^n.
\]
\label{least square method}
\end{theorem}

\begin{pfbox}
Recall that $AA^\dagger$ is the orthogonal projection onto $\Ima(A)$.
Thus
\[
AA^\dagger\mathbf{b}
=
P_{\Ima(A)}(\mathbf{b}).
\]
By the minimal distance property of orthogonal projections, for any vector
$\mathbf{y}\in \Ima(A)$, we have
\[
\|P_{\Ima(A)}(\mathbf{b})-\mathbf{b}\|
\le
\|\mathbf{y}-\mathbf{b}\|.
\]
Since any $\mathbf{y}\in \Ima(A)$ can be written as
\[
\mathbf{y}=A\mathbf{x}
\]
for some $\mathbf{x}\in \mathbb{C}^n$, we get
\[
\|A(A^\dagger\mathbf{b})-\mathbf{b}\|
\le
\|A\mathbf{x}-\mathbf{b}\|.
\]
\end{pfbox}

\begin{corollary}[Linear Regression]
For the linear regression problem, finding
$\mathbf{x}=(c_1,\dots,c_n)^t$ that minimizes the sum of squared errors
$\|A\mathbf{x}-\mathbf{b}\|^2$
is achieved by $\mathbf{x}=A^\dagger\mathbf{b}$.
\end{corollary}

\begin{theorem}[Best Solution Problem]
Let $A\in M_{m\times n}(\mathbb{C})$ and fix $\mathbf{b}\in\mathbb{C}^m$.
Suppose the system of equations $A\mathbf{x}=\mathbf{b}$ is inconsistent. Then
$\mathbf{x}=A^\dagger \mathbf{b}$ gives the best approximation of a solution of minimal norm, i.e. it satisfies:
\begin{enumerate}
    \item
    \[
    \|A(A^\dagger\mathbf{b})-\mathbf{b}\|
    \le
    \|A\mathbf{x}-\mathbf{b}\|
    \qquad
    \text{for all } \mathbf{x}\in\mathbb{C}^n.
    \]
    So $A^\dagger\mathbf{b}$ is a best solution.

    \item
    If $\mathbf{z}\in\mathbb{C}^n$ satisfies
    \[
    \|A\mathbf{z}-\mathbf{b}\|
    =
    \|A(A^\dagger\mathbf{b})-\mathbf{b}\|,
    \]
    then
    \[
    \|A^\dagger\mathbf{b}\|\le \|\mathbf{z}\|.
    \]
    So $A^\dagger\mathbf{b}$ has the minimal norm among all the best solutions.
\end{enumerate}
\end{theorem}

\begin{pfbox}
(1) has been proved in theorem~\eqref{least square method}.
For (2), suppose $\mathbf{z}$ is another best solution, $\|A\mathbf{z}-\mathbf{b}\|=\|A(A^\dagger\mathbf{b})-\mathbf{b}\|$.
Recall that if $E$ is a subspace and $\mathbf{v}\in V$, then
\[
\|P_E(\mathbf{v})-\mathbf{v}\|
\le
\|\mathbf{z}-\mathbf{v}\| \quad \forall \mathbf{z}\in E.
\]
Indeed,
\[
\mathbf{z}-\mathbf{v}
=
\underbrace{\bigl(P_E(\mathbf{v})-\mathbf{v}\bigr)}_{\in E^\perp}
+
\underbrace{\bigl(\mathbf{z}-P_E(\mathbf{v})\bigr)}_{\in E}.
\]
Hence, by the Pythagorean theorem,
\[
\|\mathbf{z}-\mathbf{v}\|^2
=
\|P_E(\mathbf{v})-\mathbf{v}\|^2
+
\|\mathbf{z}-P_E(\mathbf{v})\|^2
\ge
\|P_E(\mathbf{v})-\mathbf{v}\|^2.
\]
The equality holds when $\mathbf{z}=P_E(\mathbf{v})$.

Now,
\[
AA^\dagger\mathbf{b}
=
P_{\Ima(A)}\mathbf{b}
\in \Ima(A),
\qquad
A\mathbf{z}\in \Ima(A).
\]
Also,
\begin{align*}
\|A\mathbf{z}-\mathbf{b}\|^2
&=
\|P_{\Ima(A)}\mathbf{b}-A\mathbf{z}\|^2
+
\|P_{\Ima(A)}\mathbf{b}-\mathbf{b}\|^2.
\end{align*}
Since $\mathbf{z}$ is another best solution,
$\|A\mathbf{z}-\mathbf{b}\|=\|A(A^\dagger\mathbf{b})-\mathbf{b}\|$,
\[
A\mathbf{z}=P_{\Ima(A)}\mathbf{b}=A(A^\dagger\mathbf{b}).
\]
Multiply by $A^\dagger$ and by the Penrose conditions,
\[
\underbrace{A^\dagger A\mathbf{z}}_{P_{\ker(A)^\perp}(\mathbf{z})}=A^\dagger A(A^\dagger\mathbf{b}) = A^\dagger\mathbf{b}.
\]

Thus $A^\dagger\mathbf{b}$ is the projection of $\mathbf{z}$ onto $\ker(A)^\perp$,
we have $\|A^\dagger\mathbf{b}\|\le\|\mathbf{z}\|$.
\end{pfbox}

\begin{theorem}[Minimal Solution Problem]
Suppose the system $A\mathbf{x}=\mathbf{b}$ is consistent.
$\mathbf{x}=A^\dagger\mathbf{b}$ gives the solution of minimal norm.
If $\mathbf{z}\in \mathbb{C}^n$ is any other solution, then
\[
\|A^\dagger\mathbf{b}\|
\le
\|\mathbf{z}\|.
\]
\end{theorem}

\begin{pfbox}
Given $A\mathbf{x}=\mathbf{b}$ has a solution, we have
$\mathbf{b}\in \Ima(A)$.
Then $A(A^\dagger\mathbf{b})=P_{\Ima(A)}(\mathbf{b})=\mathbf{b}$,
so $\mathbf{x}=A^\dagger\mathbf{b}$ is a solution.

Suppose $\mathbf{z}$ satisfies $A\mathbf{z}=\mathbf{b}$.
Multiply by $A^\dagger$: $A^\dagger A\mathbf{z}=A^\dagger\mathbf{b}$.
Since $A^\dagger A=P_{\Ker(A)^\perp}$,
this means $P_{\Ker(A)^\perp}(\mathbf{z})=A^\dagger\mathbf{b}$.
Because the projection of $\mathbf{z}$ is $A^\dagger\mathbf{b}$, its norm must be
less than or equal to the norm of the original vector $\mathbf{z}$.
Therefore,
\[
\|A^\dagger\mathbf{b}\|
\le
\|\mathbf{z}\|.
\]
\end{pfbox}

\begin{remark}
    The power of pseudoinverse is that the method works in 'rank-deficient' setting.
\end{remark}




\subsection{Conclusion: SVD knows everything}

\begin{theorem}
Let $A\in M_{m\times n}(\mathbb{C})$, and suppose that a singular value decomposition of $A$ is given by
\[
A
=
U
\begin{pmatrix}
\Sigma & \mathbf{O}\\
\mathbf{O} & \mathbf{O}
\end{pmatrix}
V^*
\]
with
\[
\sigma_1\ge \sigma_2\ge \dots \ge \sigma_k>0.
\]
Let
\[
\{\mathbf{u}_1,\dots,\mathbf{u}_m\}
\]
be the columns of $U$ and
\[
\{\mathbf{v}_1,\dots,\mathbf{v}_n\}
\]
be the columns of $V$. Then:
\begin{enumerate}
    \item[(a)] $\rank(A)=k$, $\nullity(A)=n-k$.
    \item[(b)] A basis of $\Ker(A)$ is given by $\{\mathbf{v}_{k+1},\dots,\mathbf{v}_n\}$.
    \item[(c)] A basis of $\Ima(A)$ is given by $\{\mathbf{u}_1,\dots,\mathbf{u}_k\}$.
    \item[(d)] A basis of $\Ker(A^*)$ is given by $\{\mathbf{u}_{k+1},\dots,\mathbf{u}_m\}$.
    \item[(e)] A basis of $\Ima(A^*)$ is given by $\{\mathbf{v}_1,\dots,\mathbf{v}_k\}$.
    \item[(f)] The vector form of the SVD of $A$ is
    \[
    A=\sum_{i=1}^k \sigma_i\mathbf{u}_i\mathbf{v}_i^*.
    \]
    \item[(g)] The pseudoinverse of $A$ is
    \[
    A^\dagger
    =
    V
    \begin{pmatrix}
    \Sigma^{-1} & \mathbf{O}\\
    \mathbf{O} & \mathbf{O}
    \end{pmatrix}
    U^*.
    \]
    \item[(h)] The vector form of the pseudoinverse is
    \[
    A^\dagger
    =
    \sum_{i=1}^k \sigma_i^{-1}\mathbf{v}_i\mathbf{u}_i^*.
    \]
    \item[(i)] The orthogonal projection onto $\Ima(A)$ is
    \[
    AA^\dagger
    =
    \sum_{i=1}^k \mathbf{u}_i\mathbf{u}_i^*.
    \]
    \item[(j)] The orthogonal projection onto $\Ima(A^*)$ is
    \[
    A^\dagger A
    =
    \sum_{i=1}^k \mathbf{v}_i\mathbf{v}_i^*.
    \]
    \item[(k)] The orthogonal projection onto $\Ker(A)$ is
    \[
    I_n-A^\dagger A
    =
    \sum_{i=k+1}^n \mathbf{v}_i\mathbf{v}_i^*.
    \]
    \item[(l)] The orthogonal projection onto $\Ker(A^*)$ is
    \[
    I_m-AA^\dagger
    =
    \sum_{i=k+1}^m \mathbf{u}_i\mathbf{u}_i^*.
    \]
    \item[(m)] For every $\mathbf{x}\in\mathbb{C}^n$, one has
    \[
    \norm{A\mathbf{x}}^2
    =
    \sum_{i=1}^k \sigma_i^2\abs{\innerproduct{\mathbf{x}}{\mathbf{v}_i}}^2.
    \]
    \item[(n)] The maximum value of $\norm{A\mathbf{x}}$ subject to $\mathbf{x}\in\mathbb{C}^n$ with $\norm{\mathbf{x}}=1$ is $\sigma_1$.
    \item[(o)] For each integer $l$ with $1\le l\le k$, the best approximation to $A$ of rank $l$ with respect to the Frobenius norm is
    \[
    A_l
    =
    \sum_{i=1}^l \sigma_i\mathbf{u}_i\mathbf{v}_i^*.
    \]
    Moreover,
    \[
    \norm{A-A_l}_F
    =
    (\sigma_{l+1}^2+\dots+\sigma_k^2)^{1/2}.
    \]
    \item[(p)] If $A$ is square and invertible, then $A^\dagger=A^{-1}$.
    \item[(q)] If $A\mathbf{x}=\mathbf{b}$ is consistent, then $A^\dagger\mathbf{b}$ is the unique solution of minimal norm among all solutions.
    \item[(r)] If $A\mathbf{x}=\mathbf{b}$ is inconsistent, then $A^\dagger\mathbf{b}$ is a least-squares solution, and indeed the unique least-squares solution of minimal norm.
\end{enumerate}
\end{theorem}


\begin{pfbox}
\textbf{Common facts.}
Since
\[
A=U\Sigma_0V^*
\]
and $U,V$ are invertible, we have
\[
\rank(A)=\rank(\Sigma_0)=k.
\]
Hence, by the rank-nullity theorem,
\[
\nullity(A)=n-k.
\]

Moreover, from
\[
AV=U\Sigma_0,
\]
we get
\[
A\mathbf{v}_i=
\begin{cases}
\sigma_i\mathbf{u}_i, & i\le k,\\
\mathbf{0}, & i>k.
\end{cases}
\]
Thus
\[
\Ker(A)=\spn\{\mathbf{v}_{k+1},\dots,\mathbf{v}_n\},
\qquad
\Ima(A)=\spn\{\mathbf{u}_1,\dots,\mathbf{u}_k\}.
\]

Similarly, taking adjoints gives
\[
A^*=V\Sigma_0^*U^*,
\qquad
A^*U=V\Sigma_0^*.
\]
So
\[
A^*\mathbf{u}_i=
\begin{cases}
\sigma_i\mathbf{v}_i, & i\le k,\\
\mathbf{0}, & i>k.
\end{cases}
\]
Therefore
\[
\Ker(A^*)=\spn\{\mathbf{u}_{k+1},\dots,\mathbf{u}_m\},
\qquad
\Ima(A^*)=\spn\{\mathbf{v}_1,\dots,\mathbf{v}_k\}.
\]

\begin{enumerate}
    \item[(i)]
    Since $P_{\Ima(A)}=AA^\dagger$, we have
    \begin{align*}
    P_{\Ima(A)}
    &=
    AA^\dagger \\
    &=
    \left(\sum_{i=1}^k \sigma_i\mathbf{u}_i\mathbf{v}_i^*\right)
    \left(\sum_{j=1}^k \sigma_j^{-1}\mathbf{v}_j\mathbf{u}_j^*\right) \\
    &=
    \sum_{i=1}^k \mathbf{u}_i\mathbf{u}_i^*.
    \end{align*}

    \item[(j)]
    Since $P_{\Ima(A^*)}=P_{\Ker(A)^\perp}=A^\dagger A$, we have
    \[
    P_{\Ima(A^*)}
    =
    A^\dagger A
    =
    \sum_{i=1}^k \mathbf{v}_i\mathbf{v}_i^*.
    \]

    \item[(k)]
    For $\mathbf{v}\in V$, write
    \[
    \mathbf{v}=\mathbf{v}_1+\mathbf{v}_2,
    \qquad
    \mathbf{v}_1\in\Ker(A),
    \quad
    \mathbf{v}_2\in\Ker(A)^\perp.
    \]
    Then
    \[
    (I_n-A^\dagger A)\mathbf{v}
    =
    \mathbf{v}_1+\mathbf{v}_2-\mathbf{v}_2
    =
    \mathbf{v}_1
    =
    P_{\Ker(A)}(\mathbf{v}).
    \]
    Hence
    \[
    I_n-A^\dagger A
    =
    \sum_{i=1}^n \mathbf{v}_i\mathbf{v}_i^*
    -
    \sum_{i=1}^k \mathbf{v}_i\mathbf{v}_i^*
    =
    \sum_{i=k+1}^n \mathbf{v}_i\mathbf{v}_i^*.
    \]

    \item[(m)]
    For $\mathbf{x}\in\mathbb{C}^n$,
    \begin{align*}
    A\mathbf{x}
    &=
    \left(\sum_{i=1}^k \sigma_i\mathbf{u}_i\mathbf{v}_i^*\right)\mathbf{x} \\
    &=
    \sum_{i=1}^k
    \sigma_i(\mathbf{v}_i^*\mathbf{x})\mathbf{u}_i \\
    &=
    \sum_{i=1}^k
    \sigma_i\innerproduct{\mathbf{x}}{\mathbf{v}_i}\mathbf{u}_i.
    \end{align*}
    Since $\{\mathbf{u}_i\}$ is an ONB,
    \[
    \norm{A\mathbf{x}}^2
    =
    \sum_{i=1}^k
    \sigma_i^2
    \abs{\innerproduct{\mathbf{x}}{\mathbf{v}_i}}^2.
    \]

    \item[(n)]
    Suppose $\norm{\mathbf{x}}=1$. Since
    \[
    \mathbf{x}
    =
    \sum_{i=1}^n
    \innerproduct{\mathbf{x}}{\mathbf{v}_i}\mathbf{v}_i,
    \qquad
    \norm{\mathbf{x}}^2
    =
    \sum_{i=1}^n
    \abs{\innerproduct{\mathbf{x}}{\mathbf{v}_i}}^2
    =
    1,
    \]
    we have
    \begin{align*}
    \norm{A\mathbf{x}}^2
    &=
    \sum_{i=1}^k
    \sigma_i^2
    \abs{\innerproduct{\mathbf{x}}{\mathbf{v}_i}}^2 \\
    &\le
    \sigma_1^2
    \sum_{i=1}^k
    \abs{\innerproduct{\mathbf{x}}{\mathbf{v}_i}}^2 \\
    &\le
    \sigma_1^2
    \sum_{i=1}^n
    \abs{\innerproduct{\mathbf{x}}{\mathbf{v}_i}}^2 \\
    &=
    \sigma_1^2.
    \end{align*}
    Thus
    \[
    \norm{A\mathbf{x}}\le \sigma_1.
    \]
    Equality is attained when $\mathbf{x}=\mathbf{v}_1$, since
    \[
    A\mathbf{v}_1=\sigma_1\mathbf{u}_1
    \quad\Longrightarrow\quad
    \norm{A\mathbf{v}_1}^2
    =
    \sigma_1^2\norm{\mathbf{u}_1}^2
    =
    \sigma_1^2.
    \]
\end{enumerate}

Parts (f), (g), (h), (l), (o), (p), (q), and (r) were proved earlier or follow by similar arguments, so they are left as exercises.
\end{pfbox}



\begin{remark}
For $A\in M_{m\times n}(\mathbb{R})$, $A$ can be regarded as a bounded linear map from $V$ to $W$. The operator norm is
\[
\norm{A}_{op}
=
\sup_{\norm{\mathbf{x}}=1}\norm{A\mathbf{x}}
=
\sigma_1.
\]
\end{remark}